\chapter{Thảo luận và Hạn chế}
\label{ch:thaoluan}

\section{Tính đúng đắn (sound) và đầy đủ (complete) của từng kỹ thuật}
Theo định nghĩa của học phần: \emph{sound} — không có false negative khi kết luận an toàn;
\emph{complete} — không có false positive.

\begin{table}[H]
\centering\small
\caption{Tính chất của các kỹ thuật đã dùng trong đề tài.}
\label{tab:sound}
\begin{tabularx}{\linewidth}{@{}l l l X@{}}
\toprule
Kỹ thuật & Sound? & Complete? & Ghi chú trong đề tài \\
\midrule
CTL trên Kripke & Có & Có & Chỉ đúng với \emph{mô hình}, không phải mã thật \\
BMC (chưa đạt CT) & Có trong $k$ & Không & \\
BMC (đạt CT) & Có & Có & v0: CT $=3$, v1: CT $=4$, đều đã đạt \\
Static scanner & Không & Không & Heuristic theo mẫu, trong một hàm, không phân tích alias \\
ASan / UBSan & Không & Có trên đường đã chạy & Chỉ một input độc hại cho mỗi hàm \\
Z3 (int32, bit-vector) & Có & Có & Với biểu thức đã mô hình hóa \\
Black-box fuzzing & Không & Có (lỗi tìm được luôn thật) & \\
Coverage-guided fuzzing & Không & Có & Kẹt ở điều kiện magic value \\
Concolic fuzzing & Không & Có & Chỉ tập con cú pháp; không mô hình hóa được thư viện ngoài \\
\bottomrule
\end{tabularx}
\end{table}

\section{Hạn chế}
\begin{itemize}
  \item \textbf{Khoảng cách mô hình -- cài đặt.} CTL/BMC chứng minh tính chất của \emph{mô hình
        trạng thái} do chúng tôi viết tay từ yêu cầu; nếu mô hình không phản ánh đúng hệ thống
        thanh toán thật thì kết luận ``v1 an toàn'' không tự động đúng cho hệ thống đó.
  \item \textbf{Mã sinh từ đặc tả.} Các hàm C/Python được sinh theo khuôn mẫu nên nhỏ và đơn
        giản (6--10 dòng mỗi hàm). Việc công cụ phát hiện được lỗi cài có chủ đích không chứng minh
        nó sẽ phát hiện được lỗi tinh vi hơn trong mã viết tay (vd lỗi qua alias hay qua nhiều hàm).
  \item \textbf{Static scanner} tự viết, dựa trên mẫu: không sound, không complete. Vì vậy mọi
        finding bộ nhớ đều được đối chiếu với ASan — trừ CWE-401.
  \item \textbf{Môi trường.} Trên macOS Apple Silicon, LeakSanitizer không được hỗ trợ nên
        CWE-401 chỉ có bằng chứng tĩnh; chạy pipeline trên Linux sẽ bổ sung log leak.
        cppcheck không được cài trong môi trường chạy, nên bước này được bỏ qua.
  \item \textbf{Sanitizer} chỉ chạy một input độc hại cho mỗi hàm; nó xác nhận lỗi tồn tại chứ
        không chứng minh v1 an toàn với mọi input — phần đó do Z3 (số học) và phân tích bản vá đảm
        nhận.
  \item \textbf{Fuzzing} không chứng minh sự vắng mặt của lỗi. Black-box tìm được lỗi chia cho 0
        phần lớn nhờ tập giá trị biên; concolic dừng khi gặp lời gọi thư viện ngoài.
  \item \textbf{Mô hình đơn luồng.} Hai worker được mô hình bằng cụm trạng thái khóa viết tay, chưa
        dùng tổ hợp song song (parallel composition) để sinh tự động mọi cách đan xen.
\end{itemize}

\section{Phản biện và bài học}
\paragraph{Nếu bắt đầu lại, sẽ thiết kế khác điều gì?}
Đưa idempotency (event ID của webhook), giới hạn thời gian chờ ở trạng thái timeout và trạng thái
kết thúc của hoàn tiền vào đặc tả \emph{ngay từ đầu}. Ba lỗi thiết kế P4, P5, P8 đều là ``thiếu
một trạng thái/chuyển tiếp'' — những lỗi rẻ nhất để sửa ở giai đoạn đặc tả (shift-left) và đắt
nhất nếu phát hiện khi vận hành.

\paragraph{Lỗ hổng nào nghi ngờ còn tồn tại nhưng chưa xác nhận?}
Xử lý webhook đồng thời trên nhiều đơn (race giữa hai lần gửi lại cùng một sự kiện) chưa được mô
hình hóa bằng tổ hợp song song; rò rỉ bộ nhớ trên nhánh lỗi của mã viết tay (không phải mã sinh ra)
chưa được kiểm tra động do thiếu LeakSanitizer trên macOS.

\paragraph{Bài học lớn nhất về phương pháp luận.}
Không một kỹ thuật đơn lẻ nào đủ: CTL tìm lỗi thiết kế mà không test hữu hạn nào chứng minh được
(lasso của P4), Z3 chứng minh bản vá số học cho \emph{mọi} input, còn fuzzing và sanitizer xác nhận
lỗi trên mã chạy thật. Bài học thứ hai đến từ chính quá trình làm: phiên bản đầu của bộ kiểm tra
CTL âm thầm báo ``đạt'' với mọi công thức nó không hiểu (ví dụ $\mathbf{AG}\,\neg(a \wedge b)$) —
một công cụ như vậy \emph{không sound}. Chúng tôi viết lại bộ kiểm tra, buộc công thức lỗi cú pháp
phải bị báo lỗi, và kiểm tra nó bằng 30 mô hình ngẫu nhiên đối chiếu với phép duyệt vét cạn.
